% -------------------------------------------------------------
\section{Fazit und Ausblick}
% -------------------------------------------------------------

Der Compton-Effekt markiert den Wendepunkt zwischen klassischer Physik und Quantenfeldtheorie.
Zum ersten Mal wurde sichtbar, dass Licht nicht nur Wellen-, sondern auch Impulsträger ist –
und dass Energie- und Impulserhaltung nur durch den Photonenbegriff erklärbar sind.

Doch die Compton-Formel bleibt ein Modell, das Photon und Elektron
noch als Teilchen in einem klassischen Stoß behandelt.
In Wirklichkeit handelt es sich um eine quantisierte Wechselwirkung
zweier Felder – des elektromagnetischen Feldes und des Elektronenfeldes.

\begin{DidaktikBox}[Von der Formel zum Feld]
	Die Quantenelektrodynamik beschreibt den Compton-Effekt
	nicht mehr mit einzelnen Photonenbahnen, sondern
	als Wechselwirkung zwischen Feldern.
	Die Compton-Streuung entsteht dort als Austausch \emph{virtueller Photonen}.
\end{DidaktikBox}

Der Compton-Effekt verbindet experimentelle Beobachtung,
Relativität und Quantenmechanik in einer einzigen Formel.
Er bildet damit die Brücke zur Quantenelektrodynamik –
der umfassenden Theorie des Lichts und seiner Wechselwirkung mit Materie.
